\section{Controlled atomic models with observable covariates}\label{sec:controlled-fibers}
Atomic transitions need not preclude a complete continuum certificate. They do preclude treating the displacement of an atom as a small total-variation error. This section uses a different reduction: an observed covariate labels finite controlled problems, and its label remains recoverable from the current state. The regime transition is allowed to depend on the action.

\begin{assumption}[Observable finite fibers]\label{ass:fibers}
The state is $(i,z)\in\{1,\ldots,n\}\times Z$, where $Z$ is standard Borel. At date $t$, the next state is $(j,\Phi_t(z))$ with probability $P_t^a(i,j;z)$. Each $\Phi_t:Z\to Z$ is a Borel bijection with Borel inverse, independent of $i,j,a$. Primitive rewards, costs, terminal payoffs, and regime probabilities are bounded Borel functions satisfying the earlier bounds. Initially, $z$ has law $\xi$ and the regime has a fixed conditional law $q$, independent of $z$.
\end{assumption}

Let $F_0$ be the identity and $F_t=\Phi_{t-1}\circ\cdots\circ\Phi_0$. The covariate label $\zeta=F_t^{-1}(z)$ is observed at date $t$. For fixed $\zeta$, evaluate the primitives at $(i,F_t(\zeta))$ to obtain a finite controlled model. Denote its common-policy constrained optimum by $v(\zeta,\varepsilon)$ under initial regime law $q$.

\begin{theorem}[Disintegration without hidden memory]\label{thm:fibers}
Under Assumption~\ref{ass:fibers}, for $0<\beta<1$ and $\varepsilon>0$,
\begin{equation}\label{eq:fiber-identity}
 \KR(\nu,\varepsilon)=\int_Z v(\zeta,\varepsilon)\,\xi(d\zeta).
\end{equation}
The infimum on the left ranges over all Borel common randomized Markov policies, not only policies that are constant on covariate cells. No exogeneity of the regime transition is required.
\end{theorem}
\begin{proof}
A Borel policy induces the fiber policy $p_t^\zeta(a\mid i)=p_t(a\mid i,F_t(\zeta))$. The Bellman equations agree at every date and regime. An all-state feasible policy consequently has fiber cost at least $v(\zeta,\varepsilon)$ for every $\zeta$, proving the lower inequality after integration.

For the reverse inequality, apply the rational-lattice construction of Theorem~\ref{thm:lattice} separately in each fiber. At any fixed denominator, there are finitely many candidate probability arrays. Their Bellman values, relaxed-feasibility tests, lowest-index minimizers, and repaired policies are Borel in $\zeta$. Uniform primitive bounds make the lattice and repair cost errors uniform in $\zeta$. Thus the finite construction supplies a Borel family of exactly feasible policies whose costs approach $v(\zeta,\varepsilon)$ uniformly. It also establishes measurability of $v$ as a limit of Borel lower endpoints. Lift this family by using $\zeta=F_t^{-1}(z)$ at the current date. Its continuation is common to every history reaching $(t,i,z)$; it remembers no unobserved initial condition. Integrate its uniform approximation error and let the denominator tend to infinity.
\end{proof}

\begin{proposition}[Computable covariate-cell enclosures]\label{prop:fiber-cells}
Partition $Z$ into finitely many Borel cells $B_c$ and choose a representative finite fiber for each cell. Suppose primitive errors throughout cell $c$ are bounded by $d_{r,c},d_{k,c},d_{g,c}$ and the finite regime matrices satisfy
\[
 \sup_{t,i,a,\zeta\in B_c}\sum_j
 |P_t^a(i,j;F_t(\zeta))-P^a_{c,t}(i,j)|\leq\alpha_c.
\]
Compute $a_c,b_c$ from \eqref{eq:primitiveJ}--\eqref{eq:primitiveC}, using these bounds. If $\varepsilon>2a_c$ for every cell, and finite computations supply lower endpoints $L_c^+$ at allowance $\varepsilon+2a_c$ and feasible-policy upper endpoints $U_c^-$ at allowance $\varepsilon-2a_c$, then
\begin{equation}\label{eq:fiber-cells}
 \sum_c\xi(B_c)\max\{0,L_c^+-b_c\}
 \leq \KR(\nu,\varepsilon)
 \leq\sum_c\xi(B_c)(U_c^-+b_c).
\end{equation}
If the primitive bounds and finite-solver widths tend uniformly to zero under refinement, these intervals converge to the original Borel-policy optimum at the fixed allowance $\varepsilon$.
\end{proposition}
\begin{proof}
Each pair of fibers has the same finite regime coordinates. Their matrix $\ell_1$ bound controls expectations of arbitrary bounded functions on these coordinates. Standard Bellman subtraction gives the operating and cost perturbation bounds, for every policy, and gives the same operating bound for optimal values. The representative model at the relaxed allowance supplies the lower bound for each true fiber; its tightened feasible policy lifts to a feasible policy in every fiber of that cell. Integrate these inequalities, using Theorem~\ref{thm:fibers}. The tolerance-to-cost modulus from Theorem~\ref{thm:repair}, together with the finite-solver widths, proves convergence. Nonnegativity justifies truncating each lower endpoint at zero.
\end{proof}

The matrix norm in this proposition is an $\ell_1$ norm on the \emph{finite regime coordinate}, not a claim that $\delta_{\Phi_t(z)}$ and $\delta_{\Phi_t(z')}$ are close in total variation. The covariate is never smeared. A fixed observed firm type is the special case $\Phi_t(z)=z$. Nonlinear invertible covariate dynamics are also permitted. By contrast, the original maintenance transitions move the continuous condition in an action-dependent manner and do not satisfy this assumption. The fiber theorem extends the controlled atomic model class with a genuine convergence mechanism; it does not relabel the original maintenance intervals as converged.
